\section{训练35期的候选、消融与公平预算}\label{sec:round2supp}
\subsection{全部候选及支持守卫}
固定专家\(I\)为本线\(\alpha=.3,k=10\)独立EWMA；\(S\)在完整排名组内等权平均后采用相同滤波与收缩，内部偏差边界均为\(\pm.3\)。无完整组历史时\(S=I\)。同面板独立控制把\(I\)也限定到完整组活动，区分池化与支持改变。冻结两线面板中两种独立输出精确一致，原外部独立专家与本轮固定独立专家在11536案例的最大分数差也为零，两者来源身份各自保留。

风险系数只使用已发专家的已释放误差。每活动排名阶段内先对物理案例计算\(a_j=\operatorname{mean}h^2\)与\(c_j=\operatorname{mean}(\epsilon h)\)，再按活动权重汇总。全历史风险用\(k=4/12\)；近期风险用衰减.85、\(k=8\)。向零ridge收缩先于截断，避免改变估计量。逐活动删一后重新估计系数的最小至最大跨度仅描述敏感性。

门控最近最多12个完整活动，要求至少8对、原始误差Pearson\(\rho\ge.15\)、\(\mathrm{RMS}(y_1-y_2)/(\mathrm{RMS}(y_1)+\mathrm{RMS}(y_2)+10^{-8})\le.5\)，差分后的相应比例至多.75；常量序列相关按零处理。全部阈值固定。选择块30个活动阶段中仅1个开启，诊断块全部关闭。门控前风险系数及删除范围继续保存，避免将实际零系数与开关前估计混淆。
\begin{table}[!htbp]\centering
\caption{训练内滚动选择：全部12个固定共享候选}\label{tab:round2adaptive}
\small\setlength{\tabcolsep}{4pt}
\begin{tabularx}{\linewidth}{L r r r}\toprule
候选 & 选择10期 / \% & 末10期 / \% & 末段log平方风险 \\
\midrule
独立EWMA（预定选择） & 4.977 & 5.918 & 0.010006 \\
固定1/4共享 & 5.247 & 5.954 & 0.009912 \\
固定1/2共享 & 5.600 & 6.109 & 0.010134 \\
两线完全共享 & 6.534 & 6.816 & 0.011522 \\
风险收缩 k=4 & 5.028 & 5.968 & 0.010038 \\
风险收缩 k=12 & 5.015 & 5.956 & 0.010027 \\
近期风险 k=8 & 5.013 & 6.017 & 0.010095 \\
稳健门控 k=8 & 4.998 & 5.918 & 0.010006 \\
同面板独立 & 4.977 & 5.918 & 0.010006 \\
风险规则仅改T1000 & 5.008 & 5.918 & 0.010007 \\
风险规则仅改T2000 & 4.983 & 6.017 & 0.010094 \\
相邻排名风险规则 & 4.985 & 5.925 & 0.009974 \\
\bottomrule\end{tabularx}
\par\vspace{3pt}\begin{minipage}{\linewidth}\footnotesize 两线主面板；选择块372案例，末段352案例；各10活动等权。次10期按候选顺序处理并列，选中独立EWMA。log平方损失与MAPE分别保留。\end{minipage}
\end{table}

\begin{table}[!htbp]\centering
\caption{末10期的共享系数与删除敏感性}\label{tab:round2lambda}
\small\setlength{\tabcolsep}{4pt}
\begin{tabularx}{\linewidth}{L r r r r}\toprule
规则 & 正系数/状态 & 均值 & 最大值 & 最大删除跨度 \\
\midrule
风险收缩 k=4 & 15/60 & 0.033 & 0.398 & 0.625 \\
风险收缩 k=12 & 15/60 & 0.027 & 0.321 & 0.500 \\
近期风险 k=8 & 27/60 & 0.072 & 0.583 & 0.802 \\
稳健门控 k=8 & 0/60 & 0.000 & 0.000 & 0.802 \\
风险规则仅改T1000 & 6/60 & 0.009 & 0.152 & 0.365 \\
风险规则仅改T2000 & 21/60 & 0.063 & 0.583 & 0.802 \\
相邻排名风险规则 & 19/60 & 0.265 & 1.000 & 1.000 \\
\bottomrule\end{tabularx}
\par\vspace{3pt}\begin{minipage}{\linewidth}\footnotesize 60状态为10活动×2排名×3阶段。门控行的删除跨度属于门控前风险估计；实际系数全为零。删除跨度描述有限支持敏感性。\end{minipage}
\end{table}


相邻组按19档顺序固定成对，末尾单线保持独立；未用相关矩阵重排。T1000所在(1000,1500)组因T1500全缺始终回退，T2000的(2000,3000)组有支持，故该候选在主两线上具有非对称信息条件。B共35次固定专家计算、420次候选混合，无基础重训；候选、输出与选择seal完整保留。

\subsection{全部18项结构消融}
各规则共用相同基础预测、冻结案例与共同分数边界。单线消融把另一线留在独立参考；仅差异对T2000偏差设零，与仅替换T1000保留T2000独立的控制分别定义。统一\(\alpha\)、统一收缩、完整支持及逐项取消收缩按固定配置执行。
\begin{table}[!htbp]\centering
\caption{训练末10期：18项预定结构消融}\label{tab:round2ablations}
\small\setlength{\tabcolsep}{4pt}
\begin{tabularx}{\linewidth}{L r r}\toprule
固定规则 & 两线MAPE / \% & 相对独立差 / pp \\
\midrule
未校正 & 6.744 & 0.827 \\
独立EWMA & 5.918 & 0.000 \\
原锚点差异 & 6.388 & 0.471 \\
仅替换T1000 & 6.439 & 0.522 \\
仅替换T2000慢锚点 & 5.867 & -0.051 \\
T1000作锚，仅替换T2000 & 6.551 & 0.633 \\
仅锚点 & 7.035 & 1.117 \\
仅差异 & 6.355 & 0.437 \\
同滤波坐标负控 & 5.918 & 0.000 \\
同面板独立 & 5.918 & 0.000 \\
统一衰减系数 & 5.944 & 0.027 \\
统一收缩 & 6.321 & 0.403 \\
完整配对支持 & 6.388 & 0.471 \\
锚点不收缩 & 6.424 & 0.506 \\
差异不收缩 & 6.527 & 0.609 \\
全部不收缩 & 6.554 & 0.636 \\
锚点衰减改为.3 & 6.396 & 0.479 \\
差异衰减改为.3 & 5.923 & 0.006 \\
\bottomrule\end{tabularx}
\par\vspace{3pt}\begin{minipage}{\linewidth}\footnotesize 各规则共用10期352案例。方法标识对应封存配置；仅替换某一线时，另一线保留独立校正。该表保留预定消融身份，末段未用于重选。\end{minipage}
\end{table}

均值、跨期标准差、相邻期变化和lag1逐排名阶段保留于\path{C/runs/c1_002_scores/cross_event_stability.csv}及\path{event_rank_phase_residuals.csv}。稳定性量基于末10活动，阶段行保持同活动关联。原始误差、独立校正与原配方的变化共同报告，行为原因与统计机制分开解释。

\subsection{多终点历史选择、全部种子与实际预算}
N32full与N16span具有相同首末训练期，三个终点跨度分别297、305.875、341天；N16recent跨度145、153.875、175天。前缀数随冻结缺失变化。稀疏选择同时改变活动类型与组成；跨历史组编码的类别支持和\(d=56/57\)也变化。各单元自己的训练预处理对共享和独立结构完全相同，独立网络只接收所属排名监督。

共享网络无偏置\(d\to32\)编码器加两个带偏置输出头，总参数\(32d+66\)；两独立网络各为带偏置\(d\to16\to1\)，合计同为\(32d+66\)。参数相等按实际张量逐项核对。优化器为AdamW（学习率.002、权重衰减.05），全批次、无早停；活动等权且活动内排名等权。三种子集成按最终分数算术平均，非偏差或对数预测平均。

共享前向线性MAC为\(\mathrm{epochs}\,n(32d+64)\)，独立为\(\mathrm{epochs}\,n(16d+16)\)。共享每例计算两个头，独立每例只计算自己的头。相同60轮组的标量参数更新数相同，共享MAC约2.035倍；共享30轮的MAC比独立60轮多约1.73\%，每参数更新次数减半。CPU计时只覆盖拟合循环，初始化、输入读取、预处理和序列化成本另计。
\begin{table}[!htbp]\centering
\caption{公平训练的实际更新与计算预算}\label{tab:round2budget}
\small\setlength{\tabcolsep}{4pt}
\begin{tabularx}{\linewidth}{L r r r r}\toprule
结构 & 优化器步数 & 标量参数更新 & 前向线性MAC & 拟合CPU秒 \\
\midrule
independent60 & 3240 & 3044520 & 1175604480 & 0.563 \\
shared30 & 810 & 1522260 & 1195960320 & 0.249 \\
shared60 & 1620 & 3044520 & 2391920640 & 0.517 \\
\bottomrule\end{tabularx}
\par\vspace{3pt}\begin{minipage}{\linewidth}\footnotesize MAC只计训练前向线性层；反向传播、激活、偏置、优化器及预处理另有成本。共享30轮比独立60轮多约1.73\%前向MAC，每参数更新次数减半。\end{minipage}
\end{table}

\begingroup\footnotesize\setlength{\tabcolsep}{3pt}
\begin{longtable}{r p{2.5cm} r r r r}
\caption{训练期公平对照的全部27个种子单元}\label{tab:round2fairseeds}\\
\toprule
终点 & 历史组 & 种子 & 独立60轮 & 共享60轮 & 共享30轮 \\
\midrule\endfirsthead
\toprule
终点 & 历史组 & 种子 & 独立60轮 & 共享60轮 & 共享30轮 \\
\midrule\endhead
32 & N32full & 42 & 8.743 & 9.381 & 9.256 \\
32 & N32full & 123 & 7.868 & 8.804 & 8.135 \\
32 & N32full & 2026 & 8.711 & 8.138 & 7.850 \\
32 & N16span & 42 & 8.526 & 8.047 & 9.671 \\
32 & N16span & 123 & 8.001 & 9.307 & 9.803 \\
32 & N16span & 2026 & 7.933 & 9.309 & 9.522 \\
32 & N16recent & 42 & 8.986 & 8.943 & 9.545 \\
32 & N16recent & 123 & 7.502 & 8.814 & 7.595 \\
32 & N16recent & 2026 & 8.363 & 8.009 & 8.771 \\
44 & N32full & 42 & 8.705 & 7.135 & 8.324 \\
44 & N32full & 123 & 6.440 & 7.195 & 8.389 \\
44 & N32full & 2026 & 6.060 & 7.447 & 8.863 \\
44 & N16span & 42 & 10.731 & 10.080 & 9.386 \\
44 & N16span & 123 & 10.098 & 9.878 & 10.010 \\
44 & N16span & 2026 & 10.263 & 10.437 & 11.504 \\
44 & N16recent & 42 & 8.580 & 7.826 & 8.811 \\
44 & N16recent & 123 & 8.443 & 8.019 & 9.399 \\
44 & N16recent & 2026 & 6.694 & 7.073 & 8.590 \\
56 & N32full & 42 & 14.265 & 13.227 & 16.488 \\
56 & N32full & 123 & 11.271 & 11.295 & 14.290 \\
56 & N32full & 2026 & 10.697 & 11.340 & 16.488 \\
56 & N16span & 42 & 13.683 & 14.515 & 16.299 \\
56 & N16span & 123 & 11.920 & 11.828 & 15.111 \\
56 & N16span & 2026 & 11.891 & 13.073 & 16.494 \\
56 & N16recent & 42 & 13.959 & 10.867 & 12.726 \\
56 & N16recent & 123 & 9.662 & 10.589 & 12.674 \\
56 & N16recent & 2026 & 14.750 & 12.371 & 15.090 \\
\bottomrule
\end{longtable}
单位为活动等权MAPE（\%）。27个种子单元涉及108次物理拟合；共享结构各27次，独立结构54次。种子与重复终点维持活动重叠关系。\par
\endgroup


C2按结构、组、终点及种子保存MAE/RMSE/P90与配对活动差。12期汇总的总参数匹配差，独立减共享60轮在三历史组的iid区间均跨零。三个终点每个仅4期，较长块区间标注支持不足；合并12期的长块敏感性亦保留少量长段的精度限制。九单元及三个种子用于条件比较，活动数按12期计。

\section{受限机制参数、校准模拟与真实迁移}\label{sec:round2mechanism}
\subsection{原始矩、参数支持与适配缺口}
校准主面板为\(.25\le p<.9\)的T1000/T2000活动均值，前25期定参数、末10期诊断。差分矩、PSD投影、偏度峰度、JB筛查、方差前后比及滞后关联完整保存。JB为渐近筛查，前段108个检验中10个\(p<.05\)、2个BH\(q<.05\)，后段分别5和1个；样本9至25且有时间依赖，高斯适配的判断保留检验力范围。

中阶段原始\(Q\)含负特征值，\(R\)对角比为\DRatio，显示短序列波动与受限等噪声模型的适配缺口。完整900点网格固定\(m\in\{0,.5,1\}\)等离散值，描述性\(2\Delta\mathrm{NLL}\le10\)保留17点。最优共同过程项落下边界，噪声相关落上边界；支持点\(m=0\)的解释条件包括等观测方差及独立过程创新。真实共同或私有玩家状态的唯一识别仍需更多观测和更一般模型。
\begin{table}[!htbp]\centering
\caption{受限局部水平族的参数支持范围}\label{tab:round2calibration}
\small\setlength{\tabcolsep}{4pt}
\begin{tabularx}{\linewidth}{L r r r r r}\toprule
参数 & 最优 & 支持下限 & 支持上限 & 删一下限 & 删一上限 \\
\midrule
qc & 0.000000 & 0.000000 & 0.000736 & 0.000000 & 0.000000 \\
qd & 0.005888 & 0.002944 & 0.011776 & 0.002944 & 0.005888 \\
r & 0.000736 & 0.000147 & 0.001472 & 0.000736 & 0.000736 \\
rho & 0.500000 & -0.500000 & 0.500000 & 0.500000 & 0.500000 \\
m & 0.000000 & 0.000000 & 0.000000 & 0.000000 & 0.000000 \\
q\_r\_ratio & 4.000000 & 1.000000 & 45.000000 & 2.000000 & 4.000000 \\
\bottomrule\end{tabularx}
\par\vspace{3pt}\begin{minipage}{\linewidth}\footnotesize 900点离散网格；描述性支持条件为2ΔNLL≤10，保留17点。参数限于预定等观测方差与独立过程创新族。删一范围仅使用删除后所选参数；原表附带似然元数据属于完整25期。\end{minipage}
\end{table}


删一参数确由删除后似然选出，原\path{delete_one_profile.csv}附带的nll、support及样本量字段引用完整25期网格。\path{D_delete_one_metadata_audit.json}将其身份澄清为\texttt{original\_full25\_*}；本表删一范围仅使用删除后所选参数值。H另以标量滤波复算完整网格及一个固定删除点，参数选择一致。真实运行前的块接缝修订改为重抽原活动附着的真实差分贡献，未将拼接段的跳变当成真实相邻差分。

\subsection{24格条件模拟的完整失败区域}
低、中、高三个支持锚点的\(\operatorname{tr}Q/\operatorname{tr}R\)为1/4/45，三者均\(m=0\)；中位比锚点按事前顺序选实际支持点，身份与最优点分开。每锚点固定高斯、方差规范t5、5\%方差规范污染、共同突变、差异突变、1期延迟、3期延迟及异方差八种条件。每格三种子各100路径、长度35，末10期评分；共43200条路径方法记录。格间共享底层随机数，独立单位为每格的完整路径。
\begin{table}[!htbp]\centering
\caption{真实量纲校准模拟中的完整方向计数}\label{tab:round2sim}
\small\setlength{\tabcolsep}{4pt}
\begin{tabularx}{\linewidth}{L r r}\toprule
规则 & 状态MSE减少/增加 & 模拟APE减少/增加 \\
\midrule
完全共享 & 0/24 & 0/24 \\
固定半共享 & 0/24 & 0/24 \\
近期风险 & 9/15 & 9/15 \\
校准完整协方差DLM & 22/2 & 14/10 \\
已知基准协方差oracle & 22/2 & 22/2 \\
\bottomrule\end{tabularx}
\par\vspace{3pt}\begin{minipage}{\linewidth}\footnotesize 3个参数锚点×8扰动，共24格；每格3种子各100路径，方法与格间共用底层随机数。独立Monte Carlo单位为完整路径；24格描述设计覆盖。\end{minipage}
\end{table}


完全共享与半共享的24格配对Monte Carlo区间在状态MSE及模拟APE上均为正。近期风险在三个3期延迟格均增加状态风险，在共同突变格均减少相对独立风险。固定校准DLM的潜在状态MSE和模拟APE方向在部分格不同，真实得分边界也未进入模拟APE。oracle知道基准生成\(Q,R\)和初态，仍用固定局部水平递归，没有突变预知或逐时异方差信息；其扰动格负结果一并保留。

输出\texttt{raw\_state\_mse}指外部统一偏差截断之前，I/S专家内部\(\pm.3\)截断仍保留。\texttt{clipped\_fraction}只计额外外部截断；字段原名及原值未覆盖。二轮校准锚点均\(m=0\)，首轮手工网格中的载荷.5条件承担另外的机制范围。

\subsection{模拟后追加的固定DLM真实迁移}
追加计划在看过模拟后提出，再封存固定参数与代码；无额外参数搜索。仅对中阶段末10期120案例应用前25期最优\(Q,R,\mu\)及\(P_0=.1I\)，合法完整配对历史由25增至34，每期全部时点共用期初状态。H独立标量复算固定递归、当前分数边界及三比较规则，MAPE与原输出一致。

独立法为\DTransferI\%，固定DLM为\DTransfer\%，观测平方对数损失为\DTransferLog；其余逐期对照保留于\path{D/posthoc_transfer_run_001/event_comparisons.csv}。该运行逐期预测在内存收集后整体封存，未提供逐期真实公开发布时效证据。其结论覆盖固定信息集下的事后回溯迁移，正式前瞻结果由新的发布记录评价。

\section{条件可靠性、评分与未来样本量}\label{sec:round2reliability}
\subsection{完整组分层与宽度代价}
每名义水平5220案例、878完整组、16活动；同组所有冻结前缀同时覆盖才记组成功。\path{E/reliability_002/interval_conditions.csv}给76行逐案例条件，\path{interval_full_group_conditions.csv}给保持完整组的76行。排名、阶段、实际进度、类型、支持及区间来源全部保留，未据误差筛选活动。
\begin{table}[!htbp]\centering
\caption{预定条件审计的低覆盖与回退实例}\label{tab:round2conditions}
\small\setlength{\tabcolsep}{4pt}
\begin{tabularx}{\linewidth}{L r r r r}\toprule
条件 & 名义 / \% & 活动/组数 & 联合覆盖 / \% & 95\%活动区间 \\
\midrule
T1000 & 80 & 16/47 & 65.62 & [50.00,81.25] \\
T1000 & 90 & 16/47 & 87.50 & [72.92,97.92] \\
challenge类型 & 80 & 3/168 & 59.97 & [51.85,66.67] \\
challenge类型 & 90 & 3/168 & 74.79 & [62.96,84.21] \\
跨排名回退池 & 80 & 4/12 & 100.00 & [100.00,100.00] \\
跨排名回退池 & 90 & 10/27 & 100.00 & [100.00,100.00] \\
\bottomrule\end{tabularx}
\par\vspace{3pt}\begin{minipage}{\linewidth}\footnotesize 完整76个逐案例分层与76个完整组分层保留于原生CSV。条件间活动、排名、阶段及历史支持组成共同变化。回退池高覆盖伴随较宽区间。\end{minipage}
\end{table}


全活动联合0/16与1/16的固定损失bootstrap分别退化为[0,0]及活动重抽样分布；独立同分布活动二项参考另给Clopper--Pearson区间[0,20.5907\%]及[.1581\%,30.2321\%]。该二项参考另加活动独立假设。festival与medley各1期，bootstrap为单点；全成功低支持组的单点同样反映固定样本观察。支持与排名、时间和活动组成相关，条件差异解释保留这些混杂。

跨排名\texttt{all\_event\_phase\_pool}回退的完整组覆盖为100\%，80\%水平有4期12组，90\%有10期27组；其活动等权相对宽度约113\%/163\%，自身排名阶段来源约34\%/44\%。不同水平回退数量不同，较高覆盖与宽度代价共同报告。\texttt{own\_support}=0的11案例来自1期，各层小分母维持原值。

误覆盖率\(\alpha=1-\ell\)的标准Winkler评分\citep{gneiting2007proper}为
\begin{equation}
\operatorname{IS}_{\alpha}(l,u;Y)=(u-l)+\frac2\alpha(l-Y)\mathbf1(Y<l)+\frac2\alpha(Y-u)\mathbf1(Y>u).
\end{equation}
先在活动内平均再活动等权，80\%/90\%分别为\EWinklerEighty 万/\EWinklerNinety 万EP。按实际终分归一化的评分另作描述；标准分数尺度和名义水平的越界代价各自保留。最大组误差校准与逐案例区间评分回答不同目标，本轮没有按评分重新选宽度。

\subsection{规划模型、采集预算与推断条件}
功效方差只取训练35期已发legacy\_pair减独立EWMA活动MAPE差。主19线样本标准差约.079753个百分点，两线约.717920；两线规则只影响少数排名，19线平均降低差异方差。采用双侧零差检验、非中心t和正态iid卡方上95\%\(\sigma\)，再用\((1+\rho)/(1-\rho)\)作AR(1)长期方差膨胀敏感性。
\begin{table}[!htbp]\centering
\caption{训练差异驱动的两线样本量敏感性}\label{tab:round2power}
\small\setlength{\tabcolsep}{4pt}
\begin{tabularx}{\linewidth}{L r r r}\toprule
假设备择改善 / pp & ρ=0 & ρ=.3 & ρ=.6 \\
\midrule
0.25 & 151 & 279 & 597 \\
0.5 & 40 & 72 & 151 \\
1.0 & 12 & 20 & 40 \\
\bottomrule\end{tabularx}
\par\vspace{3pt}\begin{minipage}{\linewidth}\footnotesize 双侧α=.05、功效.9、训练35期正态iid卡方区间的上限σ；非中心t基于独立正态活动差，AR(1)长期方差膨胀仅为规划近似。该表规划零差检验。\end{minipage}
\end{table}


在改善.5个百分点、功效.9及\(\rho=.3\)时，主19线模型输出\ConfirmGlobalN\ 期，两线次面板为\ConfirmPairN\ 期；20期pilot在两线该条件下的模型功效约38.09\%。这些输出基于明确规划假设。60期工程采集下限与两线n72共同驱动按10向上取整的80期计划，主19线指标及双侧\(\alpha=.05\)保持不变。超过实用改善.5个百分点的移位零假设另在\path{power_usefulness_margin_grid.csv}列出，与零差检验分别规划。

未来主HAC区间采用时间排序活动差\(d_t\)：\(\widehat\gamma_h=n^{-1}\sum_{t=h+1}^n(d_t-\bar d)(d_{t-h}-\bar d)\)，\(\widehat V=\widehat\gamma_0+2\sum_{h=1}^L(1-h/(L+1))\widehat\gamma_h\)，并报告\(\bar d\pm z_{.975}\sqrt{\widehat V/n}\)。带宽\(L=\min(n-1,\max(1,\lceil n^{1/3}\rceil))\)固定。Bartlett协方差估计与弱依赖渐近条件来自既有研究\citep{newey1987hac}；零方差轨迹标记无显著性决策。iid配对t及块长1/3/5/7的固定损失敏感性另列，未来有限样本精度仍由真实依赖决定。

原20期pilot登记完整保留。新的\path{confirmation_plan.json}加独立\path{confirmation_plan_addendum.json}共同规定80期未激活队列；单独读取原JSON中的模型n4会遗漏采集补充。主比较、主指标、差值定义及双侧检验、完整活动人口、失败回退、缺失留存和固定报告时点均在激活前封存。当前未来活动为0，未来MAPE为空。
